\section{Scriptural Date and Location}
\zlabel{triptych:concise:chronology:start}
% Generated. Do not edit.
% Deterministic projection of research/chronology.toml.
% schema: 5
% document: liturgy/roman-rite/1962/propers/temporal/61-twenty-first-after-pentecost
% generated-by: tools/tpt proper-chronology annotations
\newcommand{\chronologyannotationclaim}[6]{#6}
\newcommand{\chronologyannotationreach}[3]{}
\newcommand{\chronologyannotationgroup}[3]{#3}
\newcommand{\chronologyannotation}[1]{%
  \ifcsname triptychchronologyannotation@#1\endcsname
    \csname triptychchronologyannotation@#1\endcsname
  \else
    \PackageError{triptych}{No chronology annotation for #1}{}%
  \fi}
% introit: Introit
\expandafter\def\csname triptychchronologyannotation@introit\endcsname{%
  \chronologyannotationgroup{historical-setting}{research-pending}{\textbf{Historical setting} -- Occasion date unresolved.}%
  \space \chronologyannotationgroup{composition}{nonuniform}{\textbf{Composition} -- No single date applies across every cited locus.}%
}
% epistle: Epistle
\expandafter\def\csname triptychchronologyannotation@epistle\endcsname{%
  \chronologyannotationgroup{composition}{disputed}{\textbf{Composition} -- disputed: \chronologyannotationreach{Eph.6.10}{Eph}{true}\chronologyannotationreach{Eph.6.11}{Eph}{true}\chronologyannotationreach{Eph.6.12}{Eph}{true}\chronologyannotationreach{Eph.6.13}{Eph}{true}\chronologyannotationreach{Eph.6.14}{Eph}{true}\chronologyannotationreach{Eph.6.15}{Eph}{true}\chronologyannotationreach{Eph.6.16}{Eph}{true}\chronologyannotationreach{Eph.6.17}{Eph}{true}\chronologyannotationclaim{composition.epistle-to-the-ephesians}{composition}{catholic-traditional-v1}{disputed}{a period between 58 and 63}{A.D. 58--63}; \chronologyannotationreach{Eph.6.10}{Eph}{true}\chronologyannotationreach{Eph.6.11}{Eph}{true}\chronologyannotationreach{Eph.6.12}{Eph}{true}\chronologyannotationreach{Eph.6.13}{Eph}{true}\chronologyannotationreach{Eph.6.14}{Eph}{true}\chronologyannotationreach{Eph.6.15}{Eph}{true}\chronologyannotationreach{Eph.6.16}{Eph}{true}\chronologyannotationreach{Eph.6.17}{Eph}{true}\chronologyannotationclaim{composition.epistle-to-the-ephesians}{composition}{catholic-traditional-v1}{disputed}{(Philemon; Colossians; Ephesians; Philippians), 61}{A.D. 61}.}%
}
% gradual: Gradual
\expandafter\def\csname triptychchronologyannotation@gradual\endcsname{%
  \chronologyannotationgroup{historical-setting}{research-pending}{\textbf{Historical setting} -- Occasion date unresolved.}%
  \space \chronologyannotationgroup{composition}{preferred}{\textbf{Composition}: \chronologyannotationreach{Ps.89.1}{Ps}{true}\chronologyannotationreach{Ps.89.2}{Ps}{true}\chronologyannotationclaim{critical.psalms.latest-composition-boundary}{composition}{catholic-critical-v1}{preferred}{before the Maccabean period, around 165 B.C.; no individual psalm can be dated securely}{Before c. 165 B.C.}}%
}
% alleluia: Alleluia
\expandafter\def\csname triptychchronologyannotation@alleluia\endcsname{%
  \chronologyannotationgroup{retrospective-event}{preferred}{\textbf{Retrospective event}: Preferred: \chronologyannotationreach{Ps.113.1}{Ps.77.12-13 Ps.77.42-53 Ps.104.26-38 Ps.104.43 Ps.105.7-12 Ps.105.21-22 Ps.113.1-5 Ps.135.10-15 Ps.80.6-7 2Esd.9.9-12 Acts.7.36}{false}\chronologyannotationclaim{israel.exodus.the-exodus}{retrospective-event}{catholic-traditional-v1}{preferred}{the four hundred and eightieth year after the children of Israel came out of the land of Egypt}{The Exodus from Egypt, the four hundred and eightieth year after the children of Israel came out of the land of Egypt}; alternative: \chronologyannotationreach{Ps.113.1}{Ps.77.12-13 Ps.77.42-53 Ps.104.26-38 Ps.104.43 Ps.105.7-12 Ps.105.21-22 Ps.113.1-5 Ps.135.10-15 Ps.80.6-7 2Esd.9.9-12 Acts.7.36}{false}\chronologyannotationclaim{israel.exodus.the-exodus}{retrospective-event}{catholic-traditional-v1}{alternate}{the Septuagint gives 440 years}{The Exodus from Egypt, the Septuagint gives 440 years}.}%
  \space \chronologyannotationgroup{composition}{preferred}{\textbf{Composition}: \chronologyannotationreach{Ps.113.1}{Ps}{true}\chronologyannotationclaim{critical.psalms.latest-composition-boundary}{composition}{catholic-critical-v1}{preferred}{before the Maccabean period, around 165 B.C.; no individual psalm can be dated securely}{Before c. 165 B.C.}}%
}
% gospel: Gospel
\expandafter\def\csname triptychchronologyannotation@gospel\endcsname{%
  \chronologyannotationgroup{narrated-event}{research-pending}{\textbf{Event} -- No narrated-event date in the chronology corpus.}%
  \space \chronologyannotationgroup{composition}{disputed}{\textbf{Composition} -- disputed: \chronologyannotationreach{Matt.18.23}{Matt}{true}\chronologyannotationreach{Matt.18.24}{Matt}{true}\chronologyannotationreach{Matt.18.25}{Matt}{true}\chronologyannotationreach{Matt.18.26}{Matt}{true}\chronologyannotationreach{Matt.18.27}{Matt}{true}\chronologyannotationreach{Matt.18.28}{Matt}{true}\chronologyannotationreach{Matt.18.29}{Matt}{true}\chronologyannotationreach{Matt.18.30}{Matt}{true}\chronologyannotationreach{Matt.18.31}{Matt}{true}\chronologyannotationreach{Matt.18.32}{Matt}{true}\chronologyannotationreach{Matt.18.33}{Matt}{true}\chronologyannotationreach{Matt.18.34}{Matt}{true}\chronologyannotationreach{Matt.18.35}{Matt}{true}\chronologyannotationclaim{composition.gospel-of-matthew}{composition}{catholic-traditional-v1}{disputed}{about A.D. 38-45}{c. A.D. 38--45}; \chronologyannotationreach{Matt.18.23}{Matt}{true}\chronologyannotationreach{Matt.18.24}{Matt}{true}\chronologyannotationreach{Matt.18.25}{Matt}{true}\chronologyannotationreach{Matt.18.26}{Matt}{true}\chronologyannotationreach{Matt.18.27}{Matt}{true}\chronologyannotationreach{Matt.18.28}{Matt}{true}\chronologyannotationreach{Matt.18.29}{Matt}{true}\chronologyannotationreach{Matt.18.30}{Matt}{true}\chronologyannotationreach{Matt.18.31}{Matt}{true}\chronologyannotationreach{Matt.18.32}{Matt}{true}\chronologyannotationreach{Matt.18.33}{Matt}{true}\chronologyannotationreach{Matt.18.34}{Matt}{true}\chronologyannotationreach{Matt.18.35}{Matt}{true}\chronologyannotationclaim{composition.gospel-of-matthew}{composition}{catholic-traditional-v1}{disputed}{about the year 40-42}{c. A.D. 40--42}; \chronologyannotationreach{Matt.18.23}{Matt}{true}\chronologyannotationreach{Matt.18.24}{Matt}{true}\chronologyannotationreach{Matt.18.25}{Matt}{true}\chronologyannotationreach{Matt.18.26}{Matt}{true}\chronologyannotationreach{Matt.18.27}{Matt}{true}\chronologyannotationreach{Matt.18.28}{Matt}{true}\chronologyannotationreach{Matt.18.29}{Matt}{true}\chronologyannotationreach{Matt.18.30}{Matt}{true}\chronologyannotationreach{Matt.18.31}{Matt}{true}\chronologyannotationreach{Matt.18.32}{Matt}{true}\chronologyannotationreach{Matt.18.33}{Matt}{true}\chronologyannotationreach{Matt.18.34}{Matt}{true}\chronologyannotationreach{Matt.18.35}{Matt}{true}\chronologyannotationclaim{composition.gospel-of-matthew}{composition}{catholic-traditional-v1}{disputed}{the years 40-45}{A.D. 40--45}; \chronologyannotationreach{Matt.18.23}{Matt}{true}\chronologyannotationreach{Matt.18.24}{Matt}{true}\chronologyannotationreach{Matt.18.25}{Matt}{true}\chronologyannotationreach{Matt.18.26}{Matt}{true}\chronologyannotationreach{Matt.18.27}{Matt}{true}\chronologyannotationreach{Matt.18.28}{Matt}{true}\chronologyannotationreach{Matt.18.29}{Matt}{true}\chronologyannotationreach{Matt.18.30}{Matt}{true}\chronologyannotationreach{Matt.18.31}{Matt}{true}\chronologyannotationreach{Matt.18.32}{Matt}{true}\chronologyannotationreach{Matt.18.33}{Matt}{true}\chronologyannotationreach{Matt.18.34}{Matt}{true}\chronologyannotationreach{Matt.18.35}{Matt}{true}\chronologyannotationclaim{composition.gospel-of-matthew}{composition}{catholic-traditional-v1}{disputed}{about the year 60-68}{c. A.D. 60--68}; \chronologyannotationreach{Matt.18.23}{Matt}{true}\chronologyannotationreach{Matt.18.24}{Matt}{true}\chronologyannotationreach{Matt.18.25}{Matt}{true}\chronologyannotationreach{Matt.18.26}{Matt}{true}\chronologyannotationreach{Matt.18.27}{Matt}{true}\chronologyannotationreach{Matt.18.28}{Matt}{true}\chronologyannotationreach{Matt.18.29}{Matt}{true}\chronologyannotationreach{Matt.18.30}{Matt}{true}\chronologyannotationreach{Matt.18.31}{Matt}{true}\chronologyannotationreach{Matt.18.32}{Matt}{true}\chronologyannotationreach{Matt.18.33}{Matt}{true}\chronologyannotationreach{Matt.18.34}{Matt}{true}\chronologyannotationreach{Matt.18.35}{Matt}{true}\chronologyannotationclaim{composition.gospel-of-matthew}{composition}{catholic-traditional-v1}{disputed}{about the years 64-67}{c. A.D. 64--67}; \chronologyannotationreach{Matt.18.23}{Matt}{true}\chronologyannotationreach{Matt.18.24}{Matt}{true}\chronologyannotationreach{Matt.18.25}{Matt}{true}\chronologyannotationreach{Matt.18.26}{Matt}{true}\chronologyannotationreach{Matt.18.27}{Matt}{true}\chronologyannotationreach{Matt.18.28}{Matt}{true}\chronologyannotationreach{Matt.18.29}{Matt}{true}\chronologyannotationreach{Matt.18.30}{Matt}{true}\chronologyannotationreach{Matt.18.31}{Matt}{true}\chronologyannotationreach{Matt.18.32}{Matt}{true}\chronologyannotationreach{Matt.18.33}{Matt}{true}\chronologyannotationreach{Matt.18.34}{Matt}{true}\chronologyannotationreach{Matt.18.35}{Matt}{true}\chronologyannotationclaim{composition.gospel-of-matthew}{composition}{catholic-traditional-v1}{disputed}{about the year 50}{c. A.D. 50}.}%
}
% offertory: Offertory
\expandafter\def\csname triptychchronologyannotation@offertory\endcsname{%
  \chronologyannotationgroup{composition}{preferred}{\textbf{Composition}: \chronologyannotationreach{Job.1.1}{Job}{true}\chronologyannotationreach{Job.1.12}{Job}{true}\chronologyannotationreach{Job.2.6}{Job}{true}\chronologyannotationreach{Job.2.7}{Job}{true}\chronologyannotationclaim{critical.book-of-job}{composition}{catholic-critical-v1}{preferred}{probably sometime between the seventh and fifth centuries B.C.}{Probably sometime between the seventh and fifth centuries B.C.}}%
}
% communion: Communion
\expandafter\def\csname triptychchronologyannotation@communion\endcsname{%
  \chronologyannotationgroup{traditional-attribution}{disputed}{\textbf{Traditional attribution} -- disputed: \chronologyannotationreach{Ps.118.81}{Ps.118}{true}\chronologyannotationreach{Ps.118.84}{Ps.118}{true}\chronologyannotationreach{Ps.118.86}{Ps.118}{true}\chronologyannotationclaim{israel.monarchy.david-traditional-era}{traditional-attribution}{catholic-traditional-v1}{disputed}{reigned from 1055 to 1015 B.C.}{David (reign in the usual chronology), B.C. 1055--1015}.}%
  \space \chronologyannotationgroup{historical-setting}{research-pending}{\textbf{Historical setting} -- Occasion date unresolved.}%
  \space \chronologyannotationgroup{composition}{preferred}{\textbf{Composition}: \chronologyannotationreach{Ps.118.81}{Ps}{true}\chronologyannotationreach{Ps.118.84}{Ps}{true}\chronologyannotationreach{Ps.118.86}{Ps}{true}\chronologyannotationclaim{critical.psalms.latest-composition-boundary}{composition}{catholic-critical-v1}{preferred}{before the Maccabean period, around 165 B.C.; no individual psalm can be dated securely}{Before c. 165 B.C.}}%
}
% presentation: orientation-v1
\newcommand{\chronologyorientationrow}[2]{%
  \ifcsname triptychchronologyorientation@#1\endcsname
    \csname triptychchronologyorientation@#1\endcsname{#2}%
  \else
    \PackageError{triptych}{No chronology orientation row for #1}{}%
  \fi}
\expandafter\def\csname triptychchronologyorientation@row-18bc7e647866f638\endcsname#1{%
Introit: Esth 13:9--11 & Composition: The Book of Esther. #1 & \chronologyannotationreach{Esth.13.10}{Esth}{true}\chronologyannotationreach{Esth.13.11}{Esth}{true}\chronologyannotationreach{Esth.13.9}{Esth}{true}\chronologyannotationclaim{composition.book-of-esther}{composition}{catholic-traditional-v1}{preferred}{at the end of the reign of Xerxes I (485-465 B.C.) or during the reign of his son Artaxerxes I (465-425 B.C.)}{B.C. 465--425 (Conjectural early dating; later date possible)}\\[.25em]%
}
\expandafter\def\csname triptychchronologyorientation@row-d7ec5d94ad1d0e1d\endcsname#1{%
Offertory: Job 1:1, Job 1:12, Job 2:6-7 & Composition: The Book of Job. #1 & \chronologyannotationreach{Job.1.1}{Job}{true}\chronologyannotationreach{Job.1.12}{Job}{true}\chronologyannotationreach{Job.2.6}{Job}{true}\chronologyannotationreach{Job.2.7}{Job}{true}\chronologyannotationclaim{critical.book-of-job}{composition}{catholic-critical-v1}{preferred}{probably sometime between the seventh and fifth centuries B.C.}{Probably sometime between the seventh and fifth centuries B.C.}\\[.25em]%
}
\expandafter\def\csname triptychchronologyorientation@row-b12891d28424f424\endcsname#1{%
Introit: Ps 118:1; Gradual: Psalm 89:1-2; Communion: Psalm 118:81, 84, 86 & Historical setting. #1 & Occasion date unresolved\\[.25em]%
}
\expandafter\def\csname triptychchronologyorientation@row-483ed5cea1be714a\endcsname#1{%
Introit: Ps 118:1; Gradual: Psalm 89:1-2; Alleluia: Psalm 113:1; Communion: Psalm 118:81, 84, 86 & Composition: Psalms. #1 & \chronologyannotationreach{Ps.113.1}{Ps}{true}\chronologyannotationreach{Ps.118.1}{Ps}{true}\chronologyannotationreach{Ps.118.81}{Ps}{true}\chronologyannotationreach{Ps.118.84}{Ps}{true}\chronologyannotationreach{Ps.118.86}{Ps}{true}\chronologyannotationreach{Ps.89.1}{Ps}{true}\chronologyannotationreach{Ps.89.2}{Ps}{true}\chronologyannotationclaim{critical.psalms.latest-composition-boundary}{composition}{catholic-critical-v1}{preferred}{before the Maccabean period, around 165 B.C.; no individual psalm can be dated securely}{Before c. 165 B.C.}\\[.25em]%
}
\expandafter\def\csname triptychchronologyorientation@row-602ada67f8f3bb7d\endcsname#1{%
Alleluia: Psalm 113:1 & Retrospective event: The Exodus from Egypt. #1 & Preferred: \chronologyannotationreach{Ps.113.1}{Ps.77.12-13 Ps.77.42-53 Ps.104.26-38 Ps.104.43 Ps.105.7-12 Ps.105.21-22 Ps.113.1-5 Ps.135.10-15 Ps.80.6-7 2Esd.9.9-12 Acts.7.36}{false}\chronologyannotationclaim{israel.exodus.the-exodus}{retrospective-event}{catholic-traditional-v1}{preferred}{the four hundred and eightieth year after the children of Israel came out of the land of Egypt}{Temple begun in the 480th year after the Exodus (Vulgate)}; alternative: \chronologyannotationreach{Ps.113.1}{Ps.77.12-13 Ps.77.42-53 Ps.104.26-38 Ps.104.43 Ps.105.7-12 Ps.105.21-22 Ps.113.1-5 Ps.135.10-15 Ps.80.6-7 2Esd.9.9-12 Acts.7.36}{false}\chronologyannotationclaim{israel.exodus.the-exodus}{retrospective-event}{catholic-traditional-v1}{alternate}{the Septuagint gives 440 years}{Temple begun in the 440th year after the Exodus (Septuagint)}\\[.25em]%
}
\expandafter\def\csname triptychchronologyorientation@row-b5b7957d62c25282\endcsname#1{%
Introit: Ps 118:1; Communion: Psalm 118:81, 84, 86 & Traditional attribution: David's reign. #1 & Disputed: \chronologyannotationreach{Ps.118.1}{Ps.118}{true}\chronologyannotationreach{Ps.118.81}{Ps.118}{true}\chronologyannotationreach{Ps.118.84}{Ps.118}{true}\chronologyannotationreach{Ps.118.86}{Ps.118}{true}\chronologyannotationclaim{israel.monarchy.david-traditional-era}{traditional-attribution}{catholic-traditional-v1}{disputed}{reigned from 1055 to 1015 B.C.}{B.C. 1055--1015}\\[.25em]%
}
\expandafter\def\csname triptychchronologyorientation@row-88321469d09df574\endcsname#1{%
Gospel: Matthew 18:23-35 & Event. #1 & Date unresolved\\[.25em]%
}
\expandafter\def\csname triptychchronologyorientation@row-02ce0f03aad7eafe\endcsname#1{%
Gospel: Matthew 18:23-35 & Composition: The Gospel of St Matthew. #1 & Disputed: \chronologyannotationreach{Matt.18.23}{Matt}{true}\chronologyannotationreach{Matt.18.24}{Matt}{true}\chronologyannotationreach{Matt.18.25}{Matt}{true}\chronologyannotationreach{Matt.18.26}{Matt}{true}\chronologyannotationreach{Matt.18.27}{Matt}{true}\chronologyannotationreach{Matt.18.28}{Matt}{true}\chronologyannotationreach{Matt.18.29}{Matt}{true}\chronologyannotationreach{Matt.18.30}{Matt}{true}\chronologyannotationreach{Matt.18.31}{Matt}{true}\chronologyannotationreach{Matt.18.32}{Matt}{true}\chronologyannotationreach{Matt.18.33}{Matt}{true}\chronologyannotationreach{Matt.18.34}{Matt}{true}\chronologyannotationreach{Matt.18.35}{Matt}{true}\chronologyannotationclaim{composition.gospel-of-matthew}{composition}{catholic-traditional-v1}{disputed}{about A.D. 38-45}{c. A.D. 38--45}; \chronologyannotationreach{Matt.18.23}{Matt}{true}\chronologyannotationreach{Matt.18.24}{Matt}{true}\chronologyannotationreach{Matt.18.25}{Matt}{true}\chronologyannotationreach{Matt.18.26}{Matt}{true}\chronologyannotationreach{Matt.18.27}{Matt}{true}\chronologyannotationreach{Matt.18.28}{Matt}{true}\chronologyannotationreach{Matt.18.29}{Matt}{true}\chronologyannotationreach{Matt.18.30}{Matt}{true}\chronologyannotationreach{Matt.18.31}{Matt}{true}\chronologyannotationreach{Matt.18.32}{Matt}{true}\chronologyannotationreach{Matt.18.33}{Matt}{true}\chronologyannotationreach{Matt.18.34}{Matt}{true}\chronologyannotationreach{Matt.18.35}{Matt}{true}\chronologyannotationclaim{composition.gospel-of-matthew}{composition}{catholic-traditional-v1}{disputed}{about the year 40-42}{c. A.D. 40--42}; \chronologyannotationreach{Matt.18.23}{Matt}{true}\chronologyannotationreach{Matt.18.24}{Matt}{true}\chronologyannotationreach{Matt.18.25}{Matt}{true}\chronologyannotationreach{Matt.18.26}{Matt}{true}\chronologyannotationreach{Matt.18.27}{Matt}{true}\chronologyannotationreach{Matt.18.28}{Matt}{true}\chronologyannotationreach{Matt.18.29}{Matt}{true}\chronologyannotationreach{Matt.18.30}{Matt}{true}\chronologyannotationreach{Matt.18.31}{Matt}{true}\chronologyannotationreach{Matt.18.32}{Matt}{true}\chronologyannotationreach{Matt.18.33}{Matt}{true}\chronologyannotationreach{Matt.18.34}{Matt}{true}\chronologyannotationreach{Matt.18.35}{Matt}{true}\chronologyannotationclaim{composition.gospel-of-matthew}{composition}{catholic-traditional-v1}{disputed}{the years 40-45}{A.D. 40--45}; \chronologyannotationreach{Matt.18.23}{Matt}{true}\chronologyannotationreach{Matt.18.24}{Matt}{true}\chronologyannotationreach{Matt.18.25}{Matt}{true}\chronologyannotationreach{Matt.18.26}{Matt}{true}\chronologyannotationreach{Matt.18.27}{Matt}{true}\chronologyannotationreach{Matt.18.28}{Matt}{true}\chronologyannotationreach{Matt.18.29}{Matt}{true}\chronologyannotationreach{Matt.18.30}{Matt}{true}\chronologyannotationreach{Matt.18.31}{Matt}{true}\chronologyannotationreach{Matt.18.32}{Matt}{true}\chronologyannotationreach{Matt.18.33}{Matt}{true}\chronologyannotationreach{Matt.18.34}{Matt}{true}\chronologyannotationreach{Matt.18.35}{Matt}{true}\chronologyannotationclaim{composition.gospel-of-matthew}{composition}{catholic-traditional-v1}{disputed}{about the year 60-68}{c. A.D. 60--68}; \chronologyannotationreach{Matt.18.23}{Matt}{true}\chronologyannotationreach{Matt.18.24}{Matt}{true}\chronologyannotationreach{Matt.18.25}{Matt}{true}\chronologyannotationreach{Matt.18.26}{Matt}{true}\chronologyannotationreach{Matt.18.27}{Matt}{true}\chronologyannotationreach{Matt.18.28}{Matt}{true}\chronologyannotationreach{Matt.18.29}{Matt}{true}\chronologyannotationreach{Matt.18.30}{Matt}{true}\chronologyannotationreach{Matt.18.31}{Matt}{true}\chronologyannotationreach{Matt.18.32}{Matt}{true}\chronologyannotationreach{Matt.18.33}{Matt}{true}\chronologyannotationreach{Matt.18.34}{Matt}{true}\chronologyannotationreach{Matt.18.35}{Matt}{true}\chronologyannotationclaim{composition.gospel-of-matthew}{composition}{catholic-traditional-v1}{disputed}{about the years 64-67}{c. A.D. 64--67}; \chronologyannotationreach{Matt.18.23}{Matt}{true}\chronologyannotationreach{Matt.18.24}{Matt}{true}\chronologyannotationreach{Matt.18.25}{Matt}{true}\chronologyannotationreach{Matt.18.26}{Matt}{true}\chronologyannotationreach{Matt.18.27}{Matt}{true}\chronologyannotationreach{Matt.18.28}{Matt}{true}\chronologyannotationreach{Matt.18.29}{Matt}{true}\chronologyannotationreach{Matt.18.30}{Matt}{true}\chronologyannotationreach{Matt.18.31}{Matt}{true}\chronologyannotationreach{Matt.18.32}{Matt}{true}\chronologyannotationreach{Matt.18.33}{Matt}{true}\chronologyannotationreach{Matt.18.34}{Matt}{true}\chronologyannotationreach{Matt.18.35}{Matt}{true}\chronologyannotationclaim{composition.gospel-of-matthew}{composition}{catholic-traditional-v1}{disputed}{about the year 50}{c. A.D. 50}\\[.25em]%
}
\expandafter\def\csname triptychchronologyorientation@row-d944994c29115eb5\endcsname#1{%
Epistle: Ephesians 6:10-17 & Composition: The Epistle to the Ephesians. #1 & Disputed: \chronologyannotationreach{Eph.6.10}{Eph}{true}\chronologyannotationreach{Eph.6.11}{Eph}{true}\chronologyannotationreach{Eph.6.12}{Eph}{true}\chronologyannotationreach{Eph.6.13}{Eph}{true}\chronologyannotationreach{Eph.6.14}{Eph}{true}\chronologyannotationreach{Eph.6.15}{Eph}{true}\chronologyannotationreach{Eph.6.16}{Eph}{true}\chronologyannotationreach{Eph.6.17}{Eph}{true}\chronologyannotationclaim{composition.epistle-to-the-ephesians}{composition}{catholic-traditional-v1}{disputed}{a period between 58 and 63}{A.D. 58--63}; \chronologyannotationreach{Eph.6.10}{Eph}{true}\chronologyannotationreach{Eph.6.11}{Eph}{true}\chronologyannotationreach{Eph.6.12}{Eph}{true}\chronologyannotationreach{Eph.6.13}{Eph}{true}\chronologyannotationreach{Eph.6.14}{Eph}{true}\chronologyannotationreach{Eph.6.15}{Eph}{true}\chronologyannotationreach{Eph.6.16}{Eph}{true}\chronologyannotationreach{Eph.6.17}{Eph}{true}\chronologyannotationclaim{composition.epistle-to-the-ephesians}{composition}{catholic-traditional-v1}{disputed}{(Philemon; Colossians; Ephesians; Philippians), 61}{A.D. 61}\\[.25em]%
}

\begin{chronologyorientationtable}
\chronologyorientationrow{row-18bc7e647866f638}{Story at Susa; writing-place uncertain.}
\addlinespace[2pt]
\chronologyorientationrow{row-d7ec5d94ad1d0e1d}{Story in Hus (Uz); writing-place unknown.}
\addlinespace[2pt]
\chronologyorientationrow{row-b12891d28424f424}{Occasion and site unknown.}
\addlinespace[2pt]
\chronologyorientationrow{row-483ed5cea1be714a}{Collection-wide bound; individual places unknown.}
\addlinespace[2pt]
\chronologyorientationrow{row-602ada67f8f3bb7d}{Israel leaves Egypt; interval ends at Jerusalem's Temple.}
\addlinespace[2pt]
\chronologyorientationrow{row-b5b7957d62c25282}{Conditional Davidic attribution; Psalm's setting unknown.}
\addlinespace[2pt]
\chronologyorientationrow{row-88321469d09df574}{Jesus teaches in Galilee; precise venue unknown.}
\addlinespace[2pt]
\chronologyorientationrow{row-02ce0f03aad7eafe}{Writing-place uncertain; some estimates concern a Semitic original.}
\addlinespace[2pt]
\chronologyorientationrow{row-d944994c29115eb5}{Captivity: Rome favored, Caesarea possible; readers in Asia Minor.}
\end{chronologyorientationtable}

{\small The expansive study's \textit{Historical appendix: Scriptural Date and Location} gives the full evidence: Esther, p.~23; Job and Psalms, p.~24; Matthew, p.~24; Ephesians, p.~25. Each subsection is named for its biblical book.\par
Esther: \work{Catholic Encyclopedia}, ``Esther,'' composition discussion; Esth.~1:2;2:5. Job: USCCB introduction; Job~1:1. Psalms: USCCB introduction; Haydock, Ps.~118:1, p.~781; \work{Catholic Encyclopedia}, ``David'' and ``Biblical Chronology,'' Exodus-to-Temple discussion. Matthew: Matt.~17:21--19:1; \work{Catholic Encyclopedia}, ``St. Matthew'' and ``New Testament.'' Ephesians: Eph.~6:20; \work{Catholic Encyclopedia}, ``Epistle to the Ephesians,'' date/place and recipients, and ``St. Paul,'' chronology.\par
The Esther span concerns a conjectural early date of the book, with later composition possible; it does not independently date Mordecai's Greek prayer. David's reign dates the traditional author, not this Psalm's occasion. Matthew's alternatives are historical proposals, not a single date for the extant Greek Gospel or the parable's delivery. The Psalm bound is collection-wide; the Exodus intervals date the remembered event relative to the Temple.\par}
\zlabel{triptych:concise:chronology:end}
\clearpage
